\scanpage{22}
\uline{B)} Considérons ensuite les sous-groupes \inserted{de $G$} d'ordre multiple de $3$, donc d'ordre $3$, $6$ ou $12$.

g) Sous-groupes d'ordre $3$, ce sont les $4$ $3$-sous-groupes de Sylow, indexés par les éléments de $I$
\margin{$I\quad4$}

h) Les sous-groupes d'ordre $6$ ne sont pas cycliques (il n'y a pas d'élément d'ordre $6$) donc sont isomorphes à $D_3$, \deleted{\unsure{ce}} ce sont les quatre stabilisateurs des pts de $I$, i.e. les normalisateurs des $3$-sous-groupes de Sylow
\margin{$I\quad4$}

i) Il y a un seul sous-groupe d'ordre $12$ i.e. d'indice $2$, du seul fait que $\deleted{\unsure{G_2|\mathbb{F}_2}}\unsure{\mathcal{G}_{\mathrm{ab}}|\mathfrak{S}_{\mathrm{ab}}}\simeq\{\pm1\}$, et c'est $G^+\simeq\unsure{\mathcal{G}_I^+|\mathfrak{S}_I^+}$, ss-groupe \uline{alterné} de $\unsure{\mathcal{G}_I|\mathfrak{S}_I}$
\margin{$1$}

Il y a donc \inserted{exactement} $19$ sous-groupes qui sont des deux-sous-groupes, $\ne e$, se répartissant en \uline{six} classes de conjugaison (\deleted{\unsure{bien}}\inserted{correspondant à} $6$ classes d'isomorphismes d'espaces homogènes sous $G$), et $9$ sous-groupes \inserted{$\ne G$} qui sont d'ordre multiple de $3$, se répartissant en trois autres classes de conjugaison. Cela fait en tout $28$ sous-groupes, donc $30$ en comptant $\{e\}$ et $G$, se répartissant en $9$ classes de conjugaison (comp. $11$ en comptant $\{e\}$ et $G$), qui correspondent à autant de classes d'isomorphie d'espaces homogènes sous $G$.
\margin{Pour bien faire, il faudrait les inclure dans un diagramme d'inclusions, ou diag. de morphismes d'espaces homogènes.}
% archive foliation: 21
